% C2_brief.tex -- ONE source for the English and Romanian case brief of Case 2 (Lectures 5 and 11).
\input{numbers_\lang}
\renewcommand{\emqacase}{\T{Research case 2 (Lectures 5 and 11)}{Studiul de caz 2 (cursurile 5 și 11)}}
\begin{document}
\emqatitle{\T{Case 2: Recurrent networks for an explainable price-spike early warning}%
                {Studiul de caz 2: Rețele recurente pentru avertizarea timpurie explicabilă a vîrfurilor de preț}}%
          {\T{Based on}{Studiu bazat pe lucrarea} Andrei, Pele, Cramer, Amza \& Diaconu (2026), \emph{Explainable AI-based early warning system for electricity price spikes: Evidence from the Romanian market},
           PICBE 20(1), 4820--4836, \href{https://doi.org/10.2478/picbe-2026-0352}{doi:10.2478/picbe-2026-0352}}

\begin{scenario}
\T{An industrial consumer in Romania (a steel mill, a cement plant, a data centre) can move part of its load out of an hour if it is warned a day in advance; a battery operator can keep its energy for the evening.
Every morning, before the day-ahead auction closes at noon, the energy manager asks: \textbf{will tomorrow have price spikes, in which hours, and why?}
The data team proposes a recurrent neural network that ``reads'' the last two weeks of the market hour by hour, like an experienced trader. The manager must decide \textbf{whether the network is worth it compared with the gradient-boosted trees already used in research}, and \textbf{whether its alarms can be explained} to the board before anyone acts on them.}%
{Un consumator industrial din România (o oțelărie, o fabrică de ciment, un centru de date) poate muta o parte din consum dintr-o anumită oră dacă este avertizat cu o zi înainte; un operator de baterii își poate păstra energia stocată pentru seară.
În fiecare dimineață, înainte de închiderea licitației pentru ziua următoare, la prînz, managerul de energie se întreabă: \textbf{vor exista mîine vîrfuri de preț, în ce ore și de ce?}
Echipa de analiză propune o rețea neuronală recurentă care „citește” ultimele două săptămîni ale pieței, oră cu oră, asemenea unui comerciant cu experiență. Managerul trebuie să decidă \textbf{dacă rețeaua aduce un cîștig față de arborii de decizie construiți secvențial, deja folosiți în cercetare}, și \textbf{dacă alertele ei pot fi explicate} consiliului de administrație înainte de a se acționa pe baza lor.}
\end{scenario}

\section{\T{What the paper did}{Lucrarea de referință}}
\textbf{\T{Question.}{Întrebarea de cercetare.}} \T{Can explainable artificial intelligence (XAI) identify the drivers of price spikes in the Romanian day-ahead market; is the residual load forecast a primary predictor; and can an operational early-warning system be built from interpretable balance indicators (p.~4821)?}%
{Poate inteligența artificială explicabilă (XAI, \emph{explainable artificial intelligence}) să identifice factorii vîrfurilor de preț de pe piața pentru ziua următoare din România? Este prognoza consumului rezidual una dintre principalele variabile explicative? Se poate construi un sistem operațional de avertizare timpurie pe baza unor indicatori de echilibru ușor de interpretat (p.~4821)?}

\textbf{\T{Data.}{Datele.}} \T{5,000 consecutive hours of the Romanian day-ahead market, March--September 2024, with 70 engineered features: generation, load and their day-ahead forecasts from the European Network of Transmission System Operators for Electricity (ENTSO-E), fuel and carbon prices, temperatures; a chronological 80/10/10 split that prevents data leakage (p.~4824--4826).}%
{5.000 de ore consecutive de pe piața pentru ziua următoare din România, martie--septembrie 2024, cu 70 de variabile construite: producția, consumul și prognozele lor pentru ziua următoare, publicate de Rețeaua europeană a operatorilor de transport și de sistem pentru energie electrică (ENTSO-E), prețurile combustibililor și ale certificatelor de emisii, temperaturile; o împărțire cronologică 80/10/10, care exclude scurgerea de informație (p.~4824--4826).}

\textbf{\T{Method.}{Metoda.}} \T{Gradient boosting (XGBoost, eXtreme Gradient Boosting) and Random Forest models, explained globally with SHAP (SHapley Additive exPlanations) and locally with LIME (Local Interpretable Model-agnostic Explanations);
an early-warning score that combines indicator thresholds weighted by their validation F1-score, with the alert threshold $W^*=0.477$ chosen on the validation set; a high-price event is a price above the 90th percentile, 149.99 EUR/MWh (p.~4827, 4832).}%
{Modele cu arbori de decizie construiți secvențial (XGBoost, \emph{eXtreme Gradient Boosting}) și Random Forest (păduri aleatoare), explicate global cu SHAP (\emph{SHapley Additive exPlanations}) și local cu LIME (\emph{Local Interpretable Model-agnostic Explanations});
un scor de avertizare timpurie care combină praguri ale unor indicatori, ponderate cu scorul F1 din eșantionul de validare, cu pragul de alertă $W^*=0{,}477$ ales pe eșantionul de validare; un eveniment de preț ridicat este un preț peste percentila 90, adică 149,99 EUR/MWh (p.~4827, 4832).}

\textbf{\T{Main findings of the paper.}{Principalele rezultate ale lucrării.}}
\begin{itemize}
\item \T{The best model explains 67.6\,\% of the price variance on the test set ($R^2=0.676$, mean absolute error 24.58 EUR/MWh) (p.~4828).}%
{Cel mai bun model explică 67,6\,\% din varianța prețului pe eșantionul de test ($R^2=0{,}676$, eroarea absolută medie 24,58 EUR/MWh) (p.~4828).}
\item \T{The two most influential features are the previous hour's price (mean $|$SHAP$|$ = 14.30) and the residual load forecast (10.13) (Table~5, p.~4830--4831).}%
{Cele mai influente variabile sînt prețul din ora anterioară (media $|$SHAP$|$ = 14,30) și prognoza consumului rezidual (10,13) (Tabelul~5, p.~4830--4831).}
\item \T{The warning system reaches precision 0.470, recall 0.722 and F1 0.569 on the test set (Table~7, p.~4832); the residual-load indicator alone has the highest validation F1, 0.792 (p.~4832).}%
{Pe eșantionul de test, sistemul de avertizare are precizia 0,470, sensibilitatea 0,722 și scorul F1 0,569 (Tabelul~7, p.~4832); indicatorul consumului rezidual are, singur, cel mai mare scor F1 pe eșantionul de validare, 0,792 (p.~4832).}
\item \T{Missing a spike is far more costly than responding to a false alarm, so the system favours recall (p.~4832); interpretable balance signals ``retain their directional value'' and complement the price forecasts (p.~4834).}%
{Ratarea unui vîrf costă mult mai mult decît reacția la o alarmă falsă, de aceea sistemul acordă prioritate sensibilității (p.~4832); semnalele de echilibru ușor de interpretat își păstrează valoarea ca indicatori de direcție și completează prognozele de preț (p.~4834).}
\end{itemize}
{\footnotesize\color{FooterGray}\T{Related work of our group: Andrei, Pele \& Cramer (2026), \emph{Multi-horizon explainability in energy price forecasting}, PICBE 20(1), 104--117, \href{https://doi.org/10.2478/picbe-2026-0014}{doi:10.2478/picbe-2026-0014}: on Romanian hourly data 2018--2024, short-term features (6--24 hours) carry 41\,\% of the horizon-attributed SHAP contribution and price momentum becomes 4.25 times more important in high-price regimes (abstract) --- the question of \emph{how far back} the market's memory reaches, which our recurrent network answers with integrated gradients.}%
{O lucrare înrudită a grupului nostru: Andrei, Pele și Cramer (2026), \emph{Multi-horizon explainability in energy price forecasting}, PICBE 20(1), 104--117, \href{https://doi.org/10.2478/picbe-2026-0014}{doi:10.2478/picbe-2026-0014}: pe datele orare din România, 2018--2024, variabilele pe termen scurt (6--24 de ore) concentrează 41\,\% din contribuția SHAP atribuită orizonturilor, iar inerția prețului devine de 4,25 ori mai importantă în regimurile de preț ridicat (rezumat) --- adică întrebarea \emph{cît de departe} în trecut ajunge memoria pieței, la care rețeaua noastră recurentă răspunde cu gradienții integrați.}}

\section{\T{The method, step by step}{Metoda, pas cu pas}}
\begin{intuition}
\T{Spikes happen when the system is tight, and tightness builds over days: an expensive evening yesterday, a week of weak wind, low hydro.
Gradient-boosted trees see the past only through features someone chose (``price 24 hours ago'', ``price one week ago''). A recurrent network reads the raw hourly history in order and keeps a memory of what it has seen, so it can discover such patterns itself.
To trust it, we ask which hours of the past move its alarm: integrated gradients split the network's output into contributions of every input hour, as SHAP does for the trees.}%
{Vîrfurile apar atunci cînd sistemul este sub presiune, iar presiunea se acumulează în cîteva zile: o seară scumpă ieri, o săptămînă cu vînt slab, un nivel redus al producției hidro.
Arborii de decizie construiți secvențial văd trecutul numai prin variabilele alese de analist („prețul de acum 24 de ore”, „prețul de acum o săptămînă”). O rețea recurentă citește istoricul orar brut, în ordine, și păstrează o memorie a celor văzute, astfel încît poate descoperi singură asemenea tipare.
Pentru a avea încredere în ea, ne întrebăm ce ore din trecut îi influențează alerta: gradienții integrați descompun ieșirea rețelei în contribuțiile fiecărei ore de intrare, așa cum face SHAP pentru arbori.}
\end{intuition}
\T{A spike in hour $h$ of delivery day $d$ is $S_{d,h}=\1\{p_{d,h}>\theta\}$, with $\theta$ the 90th percentile of training prices. Bids for day $d$ close at 12:00 CET (Central European Time) on day $d-1$: prices are known up to day $d-1$, load and generation up to day $d-2$.

\textbf{Paper's model.} XGBoost estimates $\hat\pi_{d,h}=\sigma\big(\sum_{k=1}^{K}f_k(\mathbf x_{d,h})\big)$, a sum of $K$ regression trees $f_k$ passed through the logistic function $\sigma(z)=1/(1+e^{-z})$; SHAP writes each log-odds prediction as $\phi_0+\sum_j\phi_{j}$, where $\phi_j$ is the Shapley value of feature $j$.
The indicator score is $W_{d,h}=\sum_m w_m\1\{x_{m,d,h}\ge\theta_m\}$, with weights $w_m\propto F1_m$; an alert is raised when $W_{d,h}\ge W^*$.

\textbf{Recurrent network (Lecture 5).} The input of day $d$ is the sequence $\mathbf u_1,\dots,\mathbf u_{T}$ of the $T=336$ hours of days $d-14,\dots,d-1$; $\mathbf u_t$ holds the standardised price of hour $t$, the load, solar, wind, residual load, hydro, gas and net imports of the same hour one day earlier, and the hour of day as $(\sin,\cos)$.
A long short-term memory (LSTM) cell updates a memory $\mathbf c_t$ and a state $\mathbf h_t$:
$\mathbf f_t=\sigma(\mathbf W_f[\mathbf h_{t-1},\mathbf u_t]+\mathbf b_f)$, $\mathbf i_t=\sigma(\mathbf W_i[\mathbf h_{t-1},\mathbf u_t]+\mathbf b_i)$, $\mathbf o_t=\sigma(\mathbf W_o[\mathbf h_{t-1},\mathbf u_t]+\mathbf b_o)$,
$\mathbf c_t=\mathbf f_t\odot\mathbf c_{t-1}+\mathbf i_t\odot\tanh(\mathbf W_c[\mathbf h_{t-1},\mathbf u_t]+\mathbf b_c)$, $\mathbf h_t=\mathbf o_t\odot\tanh(\mathbf c_t)$,
where $\mathbf f_t,\mathbf i_t,\mathbf o_t$ are the forget, input and output gates and $\odot$ is the element-wise product; the gated recurrent unit (GRU) is a lighter variant with two gates.
The final state $\mathbf h_T$ (48 units), joined with the day of week and month $\mathbf s_d$, goes through a small dense layer to 24 logits $z_{d,h}$, one per delivery hour, and $\hat\pi_{d,h}=\sigma(z_{d,h})$.
The weights minimise the binary cross-entropy $\mathcal L=-\frac1N\sum_{d,h}\big[S_{d,h}\log\hat\pi_{d,h}+(1-S_{d,h})\log(1-\hat\pi_{d,h})\big]$ with AdamW, early stopping on the validation loss, five seeds.

\textbf{Explanation over time (Lecture 11).} For an input $\mathbf u$ and the reference $\mathbf u'=\mathbf 0$ (the training mean of every channel), the integrated gradient of input $j$ (one hour, one channel) is
$\mathrm{IG}_j=(u_j-u'_j)\int_0^1\frac{\partial F(\mathbf u'+\alpha(\mathbf u-\mathbf u'))}{\partial u_j}\,d\alpha$, with $F=\sum_h z_{d,h}$ the day's total spike score; the IG sum to $F(\mathbf u)-F(\mathbf u')$, the same ``completeness'' property as SHAP.

\textbf{Quality and decision.} Precision $=TP/(TP+FP)$, recall $=TP/(TP+FN)$, $F1$ their harmonic mean, with $TP$, $FP$, $FN$ the true positives, false positives and false negatives; PR-AUC is the area under the precision--recall curve (a random rule scores the spike share); AUC is the area under the receiver operating characteristic curve.
If acting costs $C$ per MWh and an unhedged spike costs $L$ per MWh, act whenever $\hat\pi_{d,h}>C/L$.}%
{Un vîrf în ora $h$ a zilei de livrare $d$ este $S_{d,h}=\1\{p_{d,h}>\theta\}$, unde $\theta$ este percentila 90 a prețurilor din eșantionul de antrenare. Ofertele pentru ziua $d$ se închid la ora 12:00 CET (ora Europei Centrale) în ziua $d-1$: prețurile sînt cunoscute pînă în ziua $d-1$, consumul și producția pînă în ziua $d-2$.

\textbf{Modelul din lucrare.} XGBoost estimează $\hat\pi_{d,h}=\sigma\big(\sum_{k=1}^{K}f_k(\mathbf x_{d,h})\big)$, o sumă de $K$ arbori de regresie $f_k$ transformată cu funcția logistică $\sigma(z)=1/(1+e^{-z})$; SHAP descompune fiecare prognoză, pe scara logaritmului cotei de șansă, în $\phi_0+\sum_j\phi_{j}$, unde $\phi_j$ este valoarea Shapley a variabilei $j$.
Scorul bazat pe indicatori este $W_{d,h}=\sum_m w_m\1\{x_{m,d,h}\ge\theta_m\}$, cu ponderile $w_m\propto F1_m$; alerta se emite cînd $W_{d,h}\ge W^*$.

\textbf{Rețeaua recurentă (cursul 5).} Intrarea pentru ziua $d$ este secvența $\mathbf u_1,\dots,\mathbf u_{T}$ a celor $T=336$ de ore din zilele $d-14,\dots,d-1$; $\mathbf u_t$ conține prețul standardizat al orei $t$, consumul, producția solară, eoliană, consumul rezidual, producția hidro, pe gaze și importurile nete din aceeași oră a zilei precedente, precum și ora din zi, sub forma $(\sin,\cos)$.
O celulă cu memorie lungă pe termen scurt (LSTM, \emph{Long Short-Term Memory}) actualizează o memorie $\mathbf c_t$ și o stare $\mathbf h_t$:
$\mathbf f_t=\sigma(\mathbf W_f[\mathbf h_{t-1},\mathbf u_t]+\mathbf b_f)$, $\mathbf i_t=\sigma(\mathbf W_i[\mathbf h_{t-1},\mathbf u_t]+\mathbf b_i)$, $\mathbf o_t=\sigma(\mathbf W_o[\mathbf h_{t-1},\mathbf u_t]+\mathbf b_o)$,
$\mathbf c_t=\mathbf f_t\odot\mathbf c_{t-1}+\mathbf i_t\odot\tanh(\mathbf W_c[\mathbf h_{t-1},\mathbf u_t]+\mathbf b_c)$, $\mathbf h_t=\mathbf o_t\odot\tanh(\mathbf c_t)$,
unde $\mathbf f_t,\mathbf i_t,\mathbf o_t$ sînt porțile de uitare, de intrare și de ieșire, iar $\odot$ este produsul element cu element; unitatea recurentă cu porți (GRU, \emph{Gated Recurrent Unit}) este o variantă mai simplă, cu două porți.
Starea finală $\mathbf h_T$ (48 de unități), împreună cu ziua săptămînii și luna, $\mathbf s_d$, trece printr-un strat dens mic, care produce 24 de logiți $z_{d,h}$, cîte unul pentru fiecare oră de livrare, iar $\hat\pi_{d,h}=\sigma(z_{d,h})$.
Ponderile minimizează entropia încrucișată binară $\mathcal L=-\frac1N\sum_{d,h}\big[S_{d,h}\log\hat\pi_{d,h}+(1-S_{d,h})\log(1-\hat\pi_{d,h})\big]$, cu algoritmul AdamW, oprire timpurie după pierderea de validare și cinci inițializări aleatoare.

\textbf{Explicația în timp (cursul 11).} Pentru o intrare $\mathbf u$ și referința $\mathbf u'=\mathbf 0$ (media din antrenare a fiecărui canal), gradientul integrat al intrării $j$ (o oră, un canal) este
$\mathrm{IG}_j=(u_j-u'_j)\int_0^1\frac{\partial F(\mathbf u'+\alpha(\mathbf u-\mathbf u'))}{\partial u_j}\,d\alpha$, cu $F=\sum_h z_{d,h}$ scorul total de vîrf al zilei; suma gradienților integrați este $F(\mathbf u)-F(\mathbf u')$, aceeași proprietate de „completitudine” ca la SHAP.

\textbf{Calitatea și decizia.} Precizia $=TP/(TP+FP)$, sensibilitatea $=TP/(TP+FN)$, iar $F1$ este media lor armonică, unde $TP$, $FP$, $FN$ sînt numărul rezultatelor pozitive adevărate, pozitive false și negative false; PR-AUC este aria de sub curba precizie--sensibilitate (o regulă aleatoare obține ponderea vîrfurilor); AUC este aria de sub curba ROC (\emph{receiver operating characteristic}).
Dacă reacția costă $C$ pe MWh, iar un vîrf neacoperit costă $L$ pe MWh, reacționați de fiecare dată cînd $\hat\pi_{d,h}>C/L$.}
\begin{decision}
\T{Compare the network with the paper's trees and with a naive rule on the same hours, over several seeds, before paying for complexity; ask for an explanation over time for every alarm; set the action threshold from your own costs ($C/L$).}%
{Comparați rețeaua cu arborii din lucrare și cu o regulă naivă, pe aceleași ore și pe mai multe inițializări aleatoare, înainte de a plăti pentru complexitate; cereți pentru fiecare alertă o explicație în timp; stabiliți pragul de reacție din propriile costuri ($C/L$).}
\end{decision}

\section{\T{Replicate it yourself}{Replicare pe datele cursului}}
\T{Notebook \texttt{C2\_notebook\_EN.ipynb}: hourly Romanian day-ahead prices, load and generation by source (energy-charts.info), on a grid of CET delivery days. Training 2023--2024 (\NnTrainDays\ days), validation January--June 2025, test July 2025 -- October 2026 (\NnTestDays\ days, \NnTest\ hours scored by every model).
Spike threshold \Nthr\ EUR/MWh (P90 of training; paper: 149.99). Spike share: \NspikeTrain\,\% in training, \NspikeVal\,\% in validation, \NspikeTest\,\% in the test period.
Classroom implementation: our data contain no day-ahead forecasts of load, solar and wind, so all models use the realised values of day $d-2$, the last metered day at gate closure; no weather or fuel prices; about 20 features instead of 70; a longer sample with a different market regime. The LSTM and GRU networks (48 units, 14-day window) are the course's addition to the paper.}%
{Notebook-ul \texttt{C2\_notebook\_RO.ipynb} folosește prețurile orare pentru ziua următoare din România, consumul și producția pe surse (energy-charts.info), pe o grilă de zile de livrare CET. Antrenare 2023--2024 (\NnTrainDays\ de zile), validare ianuarie--iunie 2025, test iulie 2025 -- octombrie 2026 (\NnTestDays\ de zile, \NnTest\ de ore evaluate pentru toate modelele).
Pragul de vîrf este \Nthr\ EUR/MWh (percentila 90 din antrenare; în lucrare: 149,99). Ponderea orelor cu vîrfuri: \NspikeTrain\,\% în antrenare, \NspikeVal\,\% în validare, \NspikeTest\,\% în perioada de test.
Față de lucrare, implementarea de la curs ține seama de faptul că datele noastre nu conțin prognoze pentru ziua următoare ale consumului și ale producției solare și eoliene: toate modelele folosesc valorile realizate în ziua $d-2$, ultima zi măsurată la închiderea porții; nu include variabile meteorologice sau prețurile combustibililor; are aproximativ 20 de variabile în loc de 70 și un eșantion mai lung, cu un alt regim de piață. Rețelele LSTM și GRU (48 de unități, fereastră de 14 zile) sînt contribuția cursului la lucrare.}

\begin{figure}[H]\centering
\includegraphics[width=0.92\linewidth]{f1_spikes}
\caption{\T{Daily maximum price and the spike threshold (left); share of spike hours by hour of day, training and test (right).}{Prețul orar maxim al fiecărei zile și pragul de vîrf (stînga); ponderea orelor cu vîrfuri, pe ore ale zilei, în antrenare și în test (dreapta).}}
\src{\T{energy-charts.info; our computation, November 2022 -- October 2026.}{energy-charts.info; calcule proprii, noiembrie 2022 -- octombrie 2026.}}
\end{figure}

\begin{table}[H]\centering\small
\caption{\T{Early warning on the test period (July 2025 -- October 2026), alert thresholds chosen on validation by F1. s.d.\ = standard deviation of single runs across 5 seeds. Best value in bold.}{Avertizarea timpurie în perioada de test (iulie 2025 -- octombrie 2026), cu pragurile de alertă alese pe validare după F1. a.s.\ = abaterea standard a rulărilor individuale, pe 5 inițializări aleatoare. Cea mai bună valoare este marcată cu aldine.}}
\tabin{warning}
\src{\T{our computation. Paper (test set, 2024): warning system precision 0.470, recall 0.722, F1 0.569 (Table~7).}{calcule proprii. Lucrarea (eșantionul de test, 2024): sistemul de avertizare are precizia 0,470, sensibilitatea 0,722 și F1 0,569 (Tabelul~7).}}
\end{table}

\T{\textbf{What our data show.} The spike regime changed: in the test period \NspikeTest\,\% of hours are spikes (\NspikeTrain\,\% in training), concentrated in the evening (hours 18--21: \NeveTest\,\% of hours in test, \NeveTrain\,\% in training).
With gate-closure information only, the paper's XGBoost ranks hours best (AUC \NxAUC, PR-AUC \NxPR); its top SHAP features are \Ntop, \Nsecond\ and \Nthird{} --- past prices dominate, as in the paper, where the previous hour's price comes first.
The recurrent networks, which receive only the raw hourly history, reach PR-AUC \NlPR\ (LSTM; single seeds \NlPRmin--\NlPRmax) and \NgPR\ (GRU; \NgPRmin--\NgPRmax).
A block bootstrap over test days gives a PR-AUC difference LSTM $-$ XGBoost of \NdPR\ (95\,\% interval \NdPRlo\ to \NdPRhi) and an F1 difference of \NdFone\ (\NdFlo\ to \NdFhi): on two years of training data, the paper's engineered features beat the network.
In this new regime, persistence is hard to beat: the naive rule ``spike on day $d-1$ at the same hour'' has the highest F1 (\NnFone) and recall (\NnRec), which is why the paper treats recall and simple, auditable signals as central to an early-warning system.}%
{\textbf{Ce arată datele.} Regimul vîrfurilor s-a schimbat: în perioada de test \NspikeTest\,\% din ore sînt vîrfuri (\NspikeTrain\,\% în antrenare), concentrate seara (orele 18--21: \NeveTest\,\% din ore în test, \NeveTrain\,\% în antrenare).
Numai cu informația de la închiderea porții, modelul XGBoost din lucrare ordonează cel mai bine orele (AUC \NxAUC, PR-AUC \NxPR); primele variabile după SHAP sînt \Ntop, \Nsecond\ și \Nthird{} --- domină prețurile anterioare, ca în lucrare, unde prețul din ora anterioară este pe primul loc.
Rețelele recurente, care primesc numai istoricul orar brut, obțin PR-AUC \NlPR\ (LSTM; rulări individuale: \NlPRmin--\NlPRmax) și \NgPR\ (GRU; \NgPRmin--\NgPRmax).
Bootstrap-ul pe blocuri de zile de test dă o diferență de PR-AUC LSTM $-$ XGBoost de \NdPR\ (interval de 95\,\%: \NdPRlo\ pînă la \NdPRhi) și o diferență de F1 de \NdFone\ (\NdFlo\ pînă la \NdFhi): pe doi ani de date de antrenare, variabilele construite în lucrare depășesc rețeaua.
În acest regim nou, persistența este greu de depășit: regula naivă a persistenței (se semnalează un vîrf la o anumită oră dacă a existat un vîrf la aceeași oră în ziua precedentă) are cel mai mare scor F1 (\NnFone) și cea mai mare sensibilitate (\NnRec), motiv pentru care lucrarea consideră sensibilitatea și semnalele simple, verificabile, esențiale pentru un sistem de avertizare timpurie.}

\begin{figure}[H]
\begin{minipage}[t]{0.49\linewidth}\centering
\includegraphics[width=\linewidth]{f2_pr_curves}
\caption{\T{Precision--recall curves on the test period.}{Curbele precizie--sensibilitate în perioada de test.}}
\src{\T{our computation; dotted line = spike share.}{calcule proprii; linia punctată = ponderea vîrfurilor.}}
\end{minipage}\hfill
\begin{minipage}[t]{0.49\linewidth}\centering
\includegraphics[width=\linewidth]{f4_shap_xgboost}
\caption{\T{Top-10 drivers of XGBoost spike alerts (SHAP).}{Primii 10 factori ai alertelor XGBoost (SHAP).}}
\src{\T{our XGBoost classifier, seed 0; 3,000 test hours.}{clasificatorul nostru XGBoost, inițializarea 0; 3.000 de ore de test.}}
\end{minipage}
\end{figure}

\begin{figure}[H]\centering
\includegraphics[width=0.92\linewidth]{f3_time_attribution}
\caption{\T{Which hours of the past drive the LSTM's alarm? Share of $|$integrated gradients$|$ by input hour (left) and by channel (right).}{Ce ore din trecut declanșează alerta rețelei LSTM? Ponderea $|$gradienților integrați$|$ pe ore de intrare (stînga) și pe canale (dreapta).}}
\src{\T{our computation; 5 LSTM seeds, all test days, 32 integration steps.}{calcule proprii; 5 rulări LSTM, toate zilele de test, 32 de pași de integrare.}}
\end{figure}

\T{\textbf{Explaining the network.} The integrated gradients show a short memory: \NigLast\,\% of the attribution falls on the last known day ($d-1$), \NigDtwo\,\% on day $d-2$ and \NigOld\,\% on the twelve days before; within day $d-1$, the evening hours 17--21 carry \NigEve\,\%. By channel, the price carries \NigPrice\,\%, followed by \NigSecond.
The explanation also tells us \emph{why} the network trails the trees: it has not learnt the weekly cycle (day $d-7$: \NigWeek\,\%), which XGBoost receives ready-made as ``price one week earlier'', its third SHAP feature. Explaining a model is how we find what it is missing.}%
{\textbf{Explicarea rețelei.} Gradienții integrați arată o memorie scurtă: \NigLast\,\% din atribuire revine ultimei zile cunoscute ($d-1$), \NigDtwo\,\% zilei $d-2$ și \NigOld\,\% celor douăsprezece zile anterioare; în ziua $d-1$, orele de seară 17--21 concentrează \NigEve\,\%. Pe canale, prețul are \NigPrice\,\%, urmat de canalul \NigSecond.
Explicația arată și \emph{de ce} rețeaua rămîne în urma arborilor: nu a învățat ciclul săptămînal (ziua $d-7$: \NigWeek\,\%), pe care XGBoost îl primește gata construit, ca „preț cu o săptămînă înainte”, a treia variabilă după SHAP. Explicarea unui model este calea prin care aflăm ce îi lipsește.}

\begin{figure}[H]\centering
\includegraphics[width=0.6\linewidth]{f5_decision}
\caption{\T{Net value of acting by alert threshold, 1 MW of flexible load, test period.}{Valoarea netă a reacției în funcție de pragul de alertă, 1 MW de consum flexibil, perioada de test.}}
\src{\T{our computation; $C$ = 30 EUR/MWh.}{calcule proprii; $C$ = 30 EUR/MWh.}}
\end{figure}

\T{\textbf{Decision number.} With an action cost $C$ = 30 EUR/MWh (an assumption) and a loss $L$ = \Nloss\ EUR/MWh (mean spike minus mean non-spike price in training), the cost--loss threshold is $C/L$ = \Npc.
Acting at that threshold with 1 MW of flexible load would have saved about \NvalX\ thousand EUR with XGBoost and \NvalL\ thousand EUR with the LSTM over the test period; the naive rule saves \NvalN\ thousand EUR.}%
{\textbf{Cifra relevantă pentru decizie.} La un cost al reacției $C$ = 30 EUR/MWh (ipoteză) și o pierdere $L$ = \Nloss\ EUR/MWh (diferența dintre prețul mediu al orelor cu vîrfuri și cel al celorlalte ore, în antrenare), pragul cost--pierdere este $C/L$ = \Npc.
Reacționînd la acest prag cu 1 MW de consum flexibil, s-ar fi economisit în perioada de test aproximativ \NvalX\ mii EUR cu XGBoost și \NvalL\ mii EUR cu rețeaua LSTM; regula naivă economisește \NvalN\ mii EUR.}

\section{\T{Discussion questions for managers}{Întrebări pentru discuție}}
\begin{enumerate}
\item \T{On our data the paper's XGBoost beats the LSTM. What would have to change (more years of data, more input channels, day-ahead forecasts) before you would try the network again?}%
{Pe datele noastre, modelul XGBoost din lucrare depășește rețeaua LSTM. Ce ar trebui să se schimbe (mai mulți ani de date, mai multe canale de intrare, prognoze pentru ziua următoare) pentru a încerca din nou rețeaua?}
\item \T{The integrated gradients say the network looks almost only at yesterday. Is that reassuring, or does it suggest that a simpler rule would do? How would you give the network the weekly cycle?}%
{Gradienții integrați arată că rețeaua se uită aproape numai la ziua de ieri. Este un rezultat liniștitor sau sugerează că ar fi suficientă o regulă mai simplă? Cum i-ați oferi rețelei ciclul săptămînal?}
\item \T{One network seed scores higher than another on the same data. How many runs would you ask for before approving a model, and how would you report the spread?}%
{O inițializare a rețelei obține un scor mai bun decît alta, pe aceleași date. Cîte rulări ați cere înainte de a aproba un model și cum ați raporta dispersia rezultatelor?}
\item \T{What are $C$ and $L$ for your company? How would the alert threshold differ for a cement plant and for a battery operator?}%
{Care sînt valorile $C$ și $L$ pentru compania dumneavoastră? Cum ar diferi pragul de alertă pentru o fabrică de ciment și pentru un operator de baterii?}
\item \T{Spikes became daily evening events in 2025--2026. Who in the company should decide when to recalibrate the threshold and retrain the models?}%
{În 2025--2026, vîrfurile au devenit evenimente zilnice, de seară. Cine din companie ar trebui să decidă cînd se recalibrează pragul și cînd se reantrenează modelele?}
\end{enumerate}

\section{\T{Limitations of the classroom replication}{Limitele replicării}}
\T{(i) All models are fitted once (training 2023--2024, validation January--June 2025) and not re-estimated during the 16-month test. (ii) The day-ahead forecasts of load, solar and wind used in the paper are replaced by the realised values of day $d-2$; no weather or fuel prices.
(iii) The networks are small (one recurrent layer of 48 units) and are not tuned; five seeds measure their training variability. (iv) Integrated gradients depend on the reference input (here the training mean) and explain the day's total spike score, not each hour separately.
(v) The cost--loss value assumes that the flexible load can always be shifted and ignores the price impact of doing so.}%
{(i) Toate modelele sînt estimate o singură dată (antrenare 2023--2024, validare ianuarie--iunie 2025) și nu sînt reestimate în cele 16 luni de test. (ii) Prognozele pentru ziua următoare ale consumului și ale producției solare și eoliene, folosite în lucrare, sînt înlocuite cu valorile realizate în ziua $d-2$; nu se folosesc variabile meteorologice sau prețurile combustibililor.
(iii) Rețelele sînt mici (un singur strat recurent cu 48 de unități) și nu sînt optimizate; cele cinci inițializări aleatoare măsoară variabilitatea antrenării. (iv) Gradienții integrați depind de intrarea de referință (aici, media din antrenare) și explică scorul total de vîrf al zilei, nu fiecare oră separat.
(v) Calculul cost--pierdere presupune că, de fiecare dată, consumul flexibil poate fi mutat și ignoră efectul acestei mutări asupra prețului.}

\newpage
\section*{\T{Case 2 at a glance}{Studiul de caz 2 pe scurt}}
\glance{\T{Will tomorrow have price spikes, why, and is a recurrent network worth it?}{Vor exista mîine vîrfuri de preț, de ce și merită o rețea recurentă?}}%
{\T{Hourly RO prices, load and generation, 2023--Oct 2026; spike = price above P90 of training.}{Prețuri orare RO, consum și producție, 2023--oct.\ 2026; vîrf = preț peste percentila 90 din antrenare.}}%
{\T{XGBoost + SHAP (paper); LSTM and GRU on 14 days of hourly history; integrated gradients; cost--loss rule.}{XGBoost + SHAP (lucrarea); LSTM și GRU pe 14 zile de istoric orar; gradienți integrați; regula cost--pierdere.}}%
{\T{PR-AUC: XGBoost \NxPR, LSTM \NlPR; \NigLast\,\% of the network's attribution on the last known day, none on the weekly cycle.}{PR-AUC: XGBoost \NxPR, LSTM \NlPR; \NigLast\,\% din atribuirea rețelei revine ultimei zile cunoscute, nimic ciclului săptămînal.}}%
{\T{Keep the paper's trees and a persistence benchmark; act when spike probability $> C/L$ (\Npc).}{Păstrați arborii din lucrare și reperul de persistență; reacționați cînd probabilitatea unui vîrf $> C/L$ (\Npc).}}
\begin{center}
\includegraphics[width=0.8\linewidth]{fig_headline_\lang}
\src{\T{our LSTM networks (5 seeds) on energy-charts.info data; integrated gradients on the test period, July 2025 -- October 2026.}{rețelele noastre LSTM (5 rulări) pe date energy-charts.info; gradienți integrați pe perioada de test, iulie 2025 -- octombrie 2026.}}
\end{center}
\begin{minipage}[t]{0.32\linewidth}\centering{\color{MainBlue}\Huge\bfseries \NxPR}\\ \small \T{PR-AUC of XGBoost (LSTM: \NlPR)}{PR-AUC al modelului XGBoost (LSTM: \NlPR)}\end{minipage}\hfill
\begin{minipage}[t]{0.32\linewidth}\centering{\color{MainBlue}\Huge\bfseries \NspikeTest\,\%}\\ \small \T{of test hours were spikes (training: \NspikeTrain\,\%)}{din orele de test au fost vîrfuri (antrenare: \NspikeTrain\,\%)}\end{minipage}\hfill
\begin{minipage}[t]{0.32\linewidth}\centering{\color{IDAred}\Huge\bfseries \NigLast\,\%}\\ \small \T{of the LSTM's attribution on the last known day}{din atribuirea rețelei LSTM revine ultimei zile cunoscute}\end{minipage}
\end{document}
